\section{Limites projetivos em álgebra homológica}
\label{section:0.13}

\subsection{A condição de Mittag-Leffler}
\label{subsection:0.13.1}

\begin{env}[13.1.1]
\label{0.13.1.1}
Seja $\cat{C}$ uma categoria abeliana na qual existam os produtos
infinitos (axioma AB~3* de T, 1.5).
Então existe o ínfimo de uma família de subobjetos de um objeto de
$\cat{C}$, e todo sistema projetivo de objetos de $\cat{C}$ possui um
limite projetivo; este último é um funtor exato à esquerda do
sistema projetivo em questão (T, 1.8).
\oldpage[0\textsubscript{III}]{65}
Seja $(A_\alpha,f_{\alpha\beta})$ um sistema projetivo de objetos de
$\cat{C}$ cujo conjunto de índices $\mathrm{I}$ está filtrado à direita.
Ponhamos
\[
  A=\varprojlim A_\alpha,
\]
e, para cada $\alpha\in\mathrm{I}$, seja
$f_\alpha:A\longrightarrow A_\alpha$ o morfismo canônico.
Para todo $\alpha\in\mathrm{I}$, os subobjetos
$f_{\alpha\beta}(A_\beta)$, para $\alpha\leq\beta$, formam uma família
filtrada decrescente de subobjetos de $A_\alpha$.
O subobjeto
\[
  A'_\alpha
  =
  \inf_{\beta\geq\alpha}f_{\alpha\beta}(A_\beta)
\]
chama-se o subobjeto das \emph{imagens universais} em $A_\alpha$.
Está claro que
\[
  f_\alpha(A)\subset A'_\alpha,
  \qquad
  f_{\alpha\beta}(A'_\beta)\subset A'_\alpha
  \quad(\alpha\leq\beta).
\]
Assim,
$(A'_\alpha,f_{\alpha\beta}|A'_\beta)$ é um sistema projetivo e
\[
  A=\varprojlim A'_\alpha.
\]
\end{env}

\begin{env}[13.1.2]
\label{0.13.1.2}
Para um sistema projetivo $(A_\alpha,f_{\alpha\beta})$ em $\cat{C}$,
chama-se \emph{condição de Mittag-Leffler} à seguinte:
\begin{itemize}
  \item[\textbf{(ML)}]
  \emph{Para todo índice $\alpha$, existe um $\beta\geq\alpha$ tal
  que, para todo $\gamma\geq\beta$,}
  \[
    f_{\alpha\gamma}(A_\gamma)
    =
    f_{\alpha\beta}(A_\beta).
  \]
\end{itemize}

Está claro que a condição (ML) é satisfeita se todos os
$f_{\alpha\beta}$ são epimorfismos.
Reciprocamente, suponhamos que (ML) seja satisfeita, e para cada
$\alpha\in\mathrm{I}$ seja $A'_\alpha$ o subobjeto das imagens
universais em $A_\alpha$.
Então, para $\alpha\leq\beta$, a restrição de
$f_{\alpha\beta}$ a $A'_\beta$ é um \emph{epimorfismo}
\[
  A'_\beta\longrightarrow A'_\alpha.
\]
Com efeito, escolhamos $\gamma\geq\beta$ tal que
\[
  f_{\beta\delta}(A_\delta)
  =
  f_{\beta\gamma}(A_\gamma)
  \qquad(\delta\geq\gamma).
\]
Então $A'_\beta=f_{\beta\gamma}(A_\gamma)$, e portanto
\[
  f_{\alpha\beta}(A'_\beta)
  =
  f_{\alpha\gamma}(A_\gamma),
  \qquad
  A'_\alpha
  =
  f_{\alpha\gamma}(A_\gamma)
  =
  f_{\alpha\beta}(A'_\beta).
\]

A condição (ML) também é satisfeita quando os objetos $A_\alpha$ são
\emph{artinianos} em $\cat{C}$, isto é, quando toda família de
subobjetos de $A_\alpha$ possui um elemento minimal.
Com efeito, um elemento minimal da família filtrada decrescente
$\bigl(f_{\alpha\beta}(A_\beta)\bigr)$ de subobjetos de $A_\alpha$
é necessariamente o menor desses subobjetos.
\end{env}

\begin{env}[13.1.3]
\label{0.13.1.3}
\emph{Observação} (13.1.3).
A condição (ML) também pode ser formulada quando $\cat{C}$ é, por
exemplo, a categoria de conjuntos.
Pode-se definir novamente o subconjunto das imagens universais em
$A_\alpha$, e as observações feitas em \sref{0.13.1.1} e
\sref{0.13.1.2} continuam válidas.
\end{env}

\subsection{A condição de Mittag-Leffler para grupos abelianos}
\label{subsection:0.13.2}

\begin{proposition}[13.2.1]
\label{0.13.2.1}
Seja
\[
  0\longrightarrow A_\alpha
  \xrightarrow{u_\alpha}B_\alpha
  \xrightarrow{v_\alpha}C_\alpha
  \longrightarrow0
\]
uma sucessão exata de sistemas projetivos de grupos abelianos, relativos
ao mesmo conjunto de índices filtrado $\mathrm{I}$.
\begin{enumerate}
  \item[{\rm(i)}]
  Se $(B_\alpha)$ satisfaz {\rm(ML)}, então também $(C_\alpha)$.
  \item[{\rm(ii)}]
  Se $(A_\alpha)$ e $(C_\alpha)$ satisfazem {\rm(ML)}, então também
  $(B_\alpha)$.
\end{enumerate}
\end{proposition}

\begin{proof}
Sejam $(f_{\alpha\beta})$, $(g_{\alpha\beta})$ e
$(h_{\alpha\beta})$ os sistemas de homomorfismos que definem
$(A_\alpha)$, $(B_\alpha)$ e $(C_\alpha)$, respectivamente.

(i) Suponhamos que
\[
  g_{\alpha\beta}(B_\beta)
  =
  g_{\alpha\lambda}(B_\lambda)
  \qquad(\lambda\geq\beta).
\]
Como $v_\beta$ e $v_\lambda$ são sobrejetivos,
\[
  h_{\alpha\beta}(C_\beta)
  =
  v_\alpha\bigl(g_{\alpha\beta}(B_\beta)\bigr)
  =
  v_\alpha\bigl(g_{\alpha\lambda}(B_\lambda)\bigr)
  =
  h_{\alpha\lambda}(C_\lambda)
\]
para $\lambda\geq\beta$.

(ii) Seja $\alpha\in\mathrm{I}$.
Escolhamos $\beta\geq\alpha$ tal que
\[
  f_{\alpha\beta}(A_\beta)
  =
  f_{\alpha\lambda}(A_\lambda)
  \qquad(\lambda\geq\beta),
\]
e escolhamos depois $\gamma\geq\beta$ tal que
\[
  h_{\beta\gamma}(C_\gamma)
  =
  h_{\beta\lambda}(C_\lambda)
  \qquad(\lambda\geq\gamma).
\]
Seja $y_\alpha\in g_{\alpha\gamma}(B_\gamma)$, escrevamos
$y_\alpha=g_{\alpha\gamma}(y_\gamma)$ com $y_\gamma\in B_\gamma$,
e ponhamos $y_\beta=g_{\beta\gamma}(y_\gamma)$.
Então
\[
  v_\beta(y_\beta)
  =
  h_{\beta\gamma}\bigl(v_\gamma(y_\gamma)\bigr).
\]
Para todo $\lambda\geq\gamma$, existe, por hipótese, um
$y_\lambda\in B_\lambda$ tal que
\[
  h_{\beta\gamma}\bigl(v_\gamma(y_\gamma)\bigr)
  =
  h_{\beta\lambda}\bigl(v_\lambda(y_\lambda)\bigr).
\]
Daí,
$v_\beta(y_\beta-g_{\beta\lambda}(y_\lambda))=0$, e portanto
existe um $x_\beta\in A_\beta$
\oldpage[0\textsubscript{III}]{66}
tal que
\[
  y_\beta
  =
  g_{\beta\lambda}(y_\lambda)+u_\beta(x_\beta).
\]
Segue-se que
\[
  y_\alpha
  =
  g_{\alpha\lambda}(y_\lambda)
  +
  u_\alpha\bigl(f_{\alpha\beta}(x_\beta)\bigr).
\]
Como $\lambda\geq\beta$, existe um $x_\lambda\in A_\lambda$ tal
que
\[
  f_{\alpha\beta}(x_\beta)=f_{\alpha\lambda}(x_\lambda).
\]
Por conseguinte,
\[
  y_\alpha
  =
  g_{\alpha\lambda}
  \bigl(y_\lambda+u_\lambda(x_\lambda)\bigr)
  \in g_{\alpha\lambda}(B_\lambda),
\]
o que demonstra a asserção.
\end{proof}

\begin{proposition}[13.2.2]
\label{0.13.2.2}
Seja $\mathrm{I}$ um conjunto ordenado filtrado que possui um subconjunto
cofinal enumerável, e seja
\[
  0\longrightarrow A_\alpha
  \xrightarrow{u_\alpha}B_\alpha
  \xrightarrow{v_\alpha}C_\alpha
  \longrightarrow0
\]
uma sucessão exata de sistemas projetivos de grupos abelianos indexados
por $\mathrm{I}$.
Se $(A_\alpha)$ satisfaz {\rm(ML)}, então a sucessão
\[
  0\longrightarrow\varprojlim A_\alpha
  \longrightarrow\varprojlim B_\alpha
  \longrightarrow\varprojlim C_\alpha
  \longrightarrow0
\]
é exata.
\end{proposition}

\begin{proof}
Resta apenas demonstrar que o homomorfismo
\[
  v=\varprojlim v_\alpha:
  \varprojlim B_\alpha\longrightarrow\varprojlim C_\alpha
\]
é sobrejetivo.
Seja $z=(z_\alpha)\in\varprojlim C_\alpha$, e ponhamos
\[
  E_\alpha=v_\alpha^{-1}(z_\alpha).
\]
Os $E_\alpha$ formam manifestamente um sistema projetivo de conjuntos não
vazios mediante as restrições dos homomorfismos
$g_{\alpha\beta}:B_\beta\longrightarrow B_\alpha$.

Demonstremos que este sistema projetivo satisfaz (ML).
Identificando $A_\alpha$ com um subgrupo de $B_\alpha$ por meio de
$u_\alpha$, escolhamos, para todo $\alpha\in\mathrm{I}$, um índice
$\beta\geq\alpha$ tal que
\[
  g_{\alpha\beta}(A_\beta)
  =
  g_{\alpha\lambda}(A_\lambda)
  \qquad(\lambda\geq\beta).
\]
Afirmamos que também
\[
  g_{\alpha\beta}(E_\beta)
  =
  g_{\alpha\lambda}(E_\lambda)
  \qquad(\lambda\geq\beta).
\]
Com efeito, tomemos $y_\lambda\in E_\lambda$ e ponhamos
\[
  y_\beta=g_{\beta\lambda}(y_\lambda),
  \qquad
  y_\alpha=g_{\alpha\lambda}(y_\lambda).
\]
Seja $y'_\alpha=g_{\alpha\beta}(y'_\beta)$ para algum
$y'_\beta\in E_\beta$.
Então $y'_\beta-y_\beta=x_\beta\in A_\beta$, e por hipótese existe
um $x_\lambda\in A_\lambda$ tal que
\[
  g_{\alpha\beta}(x_\beta)
  =
  g_{\alpha\lambda}(x_\lambda).
\]
Assim,
\[
  y'_\alpha
  =
  g_{\alpha\beta}(y_\beta)+g_{\alpha\beta}(x_\beta)
  =
  g_{\alpha\lambda}(y_\lambda+x_\lambda)
  \in g_{\alpha\lambda}(E_\lambda),
\]
o que demonstra a afirmação.

Sabe-se (Bourbaki, \emph{Top.\ g\'{e}n.}, cap.~II,
\textsection3, teorema~1) que, sob as hipóteses impostas a
$\mathrm{I}$, um sistema projetivo de conjuntos não vazios que satisfaz (ML)
possui um limite projetivo não vazio.
Existe, portanto, um ponto
$y=(y_\alpha)\in\varprojlim E_\alpha$.
Como $v_\alpha(y_\alpha)=z_\alpha$ para todo $\alpha$, tem-se
$z=v(y)$.
\end{proof}

\begin{proposition}[13.2.3]
\label{0.13.2.3}
Sob as hipóteses sobre $\mathrm{I}$ de \sref{0.13.2.2}, seja
$(\mathrm{K}^\bullet_\alpha)_{\alpha\in\mathrm{I}}$ um sistema projetivo
de complexos de grupos abelianos
\[
  \mathrm{K}^\bullet_\alpha
  =
  (\mathrm{K}^n_\alpha)_{n\in\bb{Z}},
\]
cujas diferenciais têm grau $+1$.
Para todo $n$, existe um homomorfismo canônico funtorial
\begin{equation}
\label{0.13.2.3.1}
\tag{13.2.3.1}
  h_n:
  \HH^n\bigl(\varprojlim\mathrm{K}^\bullet_\alpha\bigr)
  \longrightarrow
  \varprojlim\HH^n(\mathrm{K}^\bullet_\alpha).
\end{equation}
Se, para todo grau $n$, o sistema projetivo de grupos abelianos
$(\mathrm{K}^n_\alpha)_{\alpha\in\mathrm{I}}$ satisfaz {\rm(ML)},
então todo $h_n$ é sobrejetivo.
Se, além disso, o sistema projetivo
$(\HH^{n-1}(\mathrm{K}^\bullet_\alpha))_{\alpha\in\mathrm{I}}$
satisfaz {\rm(ML)}, então $h_n$ é bijetivo.
\end{proposition}

\begin{proof}
Ponhamos
\[
  \mathrm{K}^n=\varprojlim\mathrm{K}^n_\alpha
  \qquad(n\in\bb{Z}).
\]
Os homomorfismos $h_n$ estão definidos pelos diagramas comutativos
\[
  \xymatrix{
    \cdots\ar[r]
      & \mathrm{K}^{n-1}\ar[r]\ar[d]
      & \mathrm{K}^n\ar[r]\ar[d]
      & \mathrm{K}^{n+1}\ar[r]\ar[d]
      & \cdots\\
    \cdots\ar[r]
      & \mathrm{K}^{n-1}_\alpha\ar[r]
      & \mathrm{K}^n_\alpha\ar[r]
      & \mathrm{K}^{n+1}_\alpha\ar[r]
      & \cdots,
  }
\]
\oldpage[0\textsubscript{III}]{67}
sendo as diferenciais em $\mathrm{K}^\bullet$ os limites projetivos
das diferenciais correspondentes nos
$\mathrm{K}^\bullet_\alpha$.
Consideremos as sucessões exatas
\[
  (*_n)\qquad
  0\longrightarrow
  \mathrm{B}^n(\mathrm{K}^\bullet_\alpha)
  \longrightarrow
  \mathrm{Z}^n(\mathrm{K}^\bullet_\alpha)
  \longrightarrow
  \HH^n(\mathrm{K}^\bullet_\alpha)
  \longrightarrow0
\]
e
\[
  (**_n)\qquad
  0\longrightarrow
  \mathrm{Z}^{n-1}(\mathrm{K}^\bullet_\alpha)
  \longrightarrow
  \mathrm{K}^{n-1}_\alpha
  \longrightarrow
  \mathrm{B}^n(\mathrm{K}^\bullet_\alpha)
  \longrightarrow0.
\]
A hipótese e a Proposição~\sref{0.13.2.1}~(i) mostram que o
sistema projetivo
$(\mathrm{B}^n(\mathrm{K}^\bullet_\alpha))_{\alpha\in\mathrm{I}}$
satisfaz (ML) para todo $n$.
Segue-se de \sref{0.13.2.2} que
\[
  (***_n)\qquad
  0\longrightarrow
  \varprojlim_\alpha\mathrm{B}^n(\mathrm{K}^\bullet_\alpha)
  \longrightarrow
  \varprojlim_\alpha\mathrm{Z}^n(\mathrm{K}^\bullet_\alpha)
  \longrightarrow
  \varprojlim_\alpha\HH^n(\mathrm{K}^\bullet_\alpha)
  \longrightarrow0
\]
é exata.
Ora,
$\varprojlim_\alpha\mathrm{B}^n(\mathrm{K}^\bullet_\alpha)$
identifica-se com um subgrupo de $\mathrm{K}^{n}$ que contém
$\mathrm{B}^n(\mathrm{K}^\bullet)$, enquanto que
$\varprojlim_\alpha\mathrm{Z}^n(\mathrm{K}^\bullet_\alpha)$
identifica-se com um subgrupo de $\mathrm{Z}^n(\mathrm{K}^\bullet)$.
Por conseguinte, $h_n$ é sobrejetivo.

Suponhamos agora, além disso, que
$(\HH^{n-1}(\mathrm{K}^\bullet_\alpha))_{\alpha\in\mathrm{I}}$
satisfaz (ML).
As sucessões exatas $(*_{n-1})$ e a
Proposição~\sref{0.13.2.1}~(ii) mostram que
$(\mathrm{Z}^{n-1}(\mathrm{K}^\bullet_\alpha))_{\alpha\in\mathrm{I}}$
satisfaz (ML).
Aplicando \sref{0.13.2.2} às sucessões exatas $(**_n)$ obtém-se uma
sucessão exata
\[
  0\longrightarrow
  \varprojlim_\alpha
    \mathrm{Z}^{n-1}(\mathrm{K}^\bullet_\alpha)
  \longrightarrow
  \mathrm{K}^{n-1}
  \xrightarrow{u}
  \varprojlim_\alpha
    \mathrm{B}^n(\mathrm{K}^\bullet_\alpha)
  \longrightarrow0.
\]
Como
$\varprojlim_\alpha\mathrm{B}^n(\mathrm{K}^\bullet_\alpha)
\supset\mathrm{B}^n(\mathrm{K}^\bullet)$
e o composto de $u$ com a injeção
\[
  \varprojlim_\alpha\mathrm{B}^n(\mathrm{K}^\bullet_\alpha)
  \longrightarrow\mathrm{K}^n
\]
é a diferencial
$\mathrm{K}^{n-1}\longrightarrow\mathrm{K}^n$, a sobrejetividade de
$u$ implica
\[
  \varprojlim_\alpha\mathrm{B}^n(\mathrm{K}^\bullet_\alpha)
  =
  \mathrm{B}^n(\mathrm{K}^\bullet).
\]
Assim, $h_n$ é injetivo.
\end{proof}

\begin{env}[13.2.4]
\label{0.13.2.4}
\emph{Observações} (13.2.4).
\begin{enumerate}
  \item[{\rm(i)}]
  O argumento de \sref{0.13.2.2}
  (cf.\ Bourbaki, \emph{loc.\ cit.}) mostra que a conclusão de
  essa proposição continua válida se supusermos apenas que os
  $A_\alpha$ podem ser dotados de estruturas de espaços metrizáveis completos
  nos quais as translações são homeomorfismos, que as aplicações
  $f_{\alpha\beta}:A_\beta\longrightarrow A_\alpha$ que definem o
  sistema projetivo são uniformemente contínuas para as métricas em
  questão, e que $(A_\alpha)$ satisfaz a condição
  \begin{itemize}
    \item[\textbf{(ML')}]
    \emph{Para todo índice $\alpha$, existe um $\beta\geq\alpha$ tal
    que, para todo $\gamma\geq\beta$,
    $f_{\alpha\gamma}(A_\gamma)$ é denso em
    $f_{\alpha\beta}(A_\beta)$.}
  \end{itemize}

  Isto dá um suplemento análogo a \sref{0.13.2.3}.
  Suponhamos que $\mathrm{K}^n_\alpha=0$ para $n<0$ e todo $\alpha$;
  que
  $(\mathrm{K}^n_\alpha)_{\alpha\in\mathrm{I}}$ satisfaz (ML) para
  $n\geq0$; e que
  $A_\alpha=\HH^0(\mathrm{K}^\bullet_\alpha)$ pode ser dotado de
  estruturas de espaço métrico que tenham as propriedades recém-enunciadas.
  Então as conclusões de \sref{0.13.2.3} não mudam para
  $n\geq2$, e $h_1$ é além disso \emph{bijetivo}.
  Com efeito, o argumento de \sref{0.13.2.2} mostra também que
  $(\mathrm{B}^1(\mathrm{K}^\bullet_\alpha))_{\alpha\in\mathrm{I}}$
  satisfaz (ML), que a sucessão $(***_1)$ é exata e, finalmente,
  que
  \[
    \varprojlim_\alpha
    \mathrm{B}^1(\mathrm{K}^\bullet_\alpha)
    =
    \mathrm{B}^1(\mathrm{K}^\bullet).
  \]
  Isto demonstra, entre outras coisas, as asserções de (T, 3.10.2).

  \item[{\rm(ii)}]
  Podem-se introduzir os funtores derivados à direita
  $\varprojlim^{(n)}$ do funtor $\varprojlim$, e obter enunciados mais
  completos que os precedentes~\cite{III-28}.
\end{enumerate}
\end{env}

\subsection{Aplicação: cohomologia de um limite projetivo de feixes}
\label{subsection:0.13.3}

\begin{proposition}[13.3.1]
\label{0.13.3.1}
\oldpage[0\textsubscript{III}]{68}
Seja $X$ um espaço topológico, seja
$(\sh{F}_k)_{k\in\bb{N}}$ um sistema projetivo de feixes de
grupos abelianos sobre $X$, e ponhamos
$\sh{F}=\varprojlim_k\sh{F}_k$.
Suponhamos que sejam satisfeitas as condições seguintes:
\begin{enumerate}
  \item[{\rm(i)}]
  Existe uma base $\mathfrak{B}$ da topologia de $X$ tal que,
  para todo $U\in\mathfrak{B}$ e todo $i\geq0$, o sistema
  projetivo
  $(\HH^i(U,\sh{F}_k))_{k\in\bb{N}}$ satisfaz (ML).

  \item[{\rm(ii)}]
  Para todo $x\in X$ e todo $i>0$, tem-se
  \[
    \varinjlim_U\bigl(\varprojlim_k
    \HH^i(U,\sh{F}_k)\bigr)=0,
  \]
  onde $U$ percorre as vizinhanças de $x$ pertencentes a
  $\mathfrak{B}$.

  \item[{\rm(iii)}]
  Os homomorfismos
  $u_{hk}:\sh{F}_k\longrightarrow\sh{F}_h$ $(h\leq k)$ que definem o
  sistema projetivo $(\sh{F}_k)$ são sobrejetivos.
\end{enumerate}
Sob estas condições, para todo $i>0$, o homomorfismo canônico
\[
  h_i:\HH^i(X,\sh{F})
  \longrightarrow
  \varprojlim_k\HH^i(X,\sh{F}_k)
\]
é sobrejetivo. Se, além disso, para um valor dado de $i$ o sistema
projetivo
$(\HH^{i-1}(X,\sh{F}_k))_{k\in\bb{N}}$ satisfaz (ML), então $h_i$ é
bijetivo.
\end{proposition}

\begin{proof}
\emph{a)}
Suponhamos primeiro que os $\sh{F}_k$, assim como os núcleos
$\sh{N}_{hk}$ dos $u_{hk}$, são flácidos; mostraremos então que
a condição~{\rm(iii)} do enunciado implica
$\HH^i(X,\sh{F})=0$ para $i>0$.
Basta demonstrar que, para todo subconjunto aberto $U$ de $X$ e todo
recobrimento $\mathfrak{U}$ de $U$ por subconjuntos abertos de $U$, tem-se
\[
  \CHH^i(\mathfrak{U},\sh{F})=0
  \qquad (i>0).
\]
Com efeito, segue-se então primeiro que, para a cohomologia de \v{C}ech,
$\check{\HH}^{\,i}(U,\sh{F})=0$ para todo $i>0$, e portanto, por
(G, II, 5.9.2) aplicado ao conjunto de todos os subconjuntos abertos de
$X$, que $\HH^i(X,\sh{F})=0$ para todo $i>0$.

Como os $\sh{F}_k$ são flácidos, tem-se
$\CHH^i(\mathfrak{U},\sh{F}_k)=0$ para todo $i>0$
(G, II, 5.2.3).
Para cada $k$, consideremos o complexo
$\check{\mathrm{C}}^\bullet(\mathfrak{U},\sh{F}_k)$ de cocadeias
alternadas sobre o nervo do recobrimento $\mathfrak{U}$ (G, II, 5.1);
estes complexos formam manifestamente um sistema projetivo de complexos de
grupos abelianos.
Demonstremos que todas as aplicações
\[
  \check{\mathrm{C}}^\bullet(\mathfrak{U},\sh{F}_k)
  \longrightarrow
  \check{\mathrm{C}}^\bullet(\mathfrak{U},\sh{F}_h)
  \qquad(h\leq k)
\]
são sobrejetivas.
Por definição, basta mostrar que, para todo subconjunto aberto $V$
de $X$, a aplicação
$\Gamma(V,\sh{F}_k)\longrightarrow\Gamma(V,\sh{F}_h)$ é
  sobrejetiva. Mas a sucessão
\[
  0\longrightarrow\sh{N}_{hk}\longrightarrow\sh{F}_k
  \longrightarrow\sh{F}_h\longrightarrow0
\]
é exata por hipótese e dá, portanto, a sucessão exata de
cohomologia
\[
  \Gamma(V,\sh{F}_k)\longrightarrow\Gamma(V,\sh{F}_h)
  \longrightarrow\HH^1(V,\sh{N}_{hk})=0,
\]
pois $\sh{N}_{hk}$ é flácido.

O sistema projetivo
$(\check{\mathrm{C}}^\bullet(\mathfrak{U},\sh{F}_k))_{k\in\bb{N}}$
satisfaz, portanto, (ML), e o mesmo vale para
$(\CHH^i(\mathfrak{U},\sh{F}_k))_{k\in\bb{N}}$ para todo $i\geq0$:
para $i>0$ isto é trivial, enquanto que para $i=0$ segue-se do
precedente, pois
$\CHH^0(\mathfrak{U},\sh{F}_k)=\Gamma(U,\sh{F}_k)$
(G, II, 5.2.2).
Podemos, portanto, aplicar \sref{0.13.2.3}, que dá
\[
  \CHH^i(\mathfrak{U},\sh{F})
  =
  \varprojlim_k\CHH^i(\mathfrak{U},\sh{F}_k)
  =0
  \qquad(i>0).
\]

\emph{b)}
Passemos agora ao caso geral. Para cada $k\in\bb{N}$, consideremos a
\emph{resolução canônica}
\[
  \sh{C}^\bullet(X,\sh{F}_k)
  =
  \bigl(\sh{C}^i(X,\sh{F}_k)\bigr)_{i\geq0}
\]
de $\sh{F}_k$ por feixes flácidos (G, II, 4.3).
Para todo $i\geq0$, está claro que
$(\sh{C}^i(X,\sh{F}_k))_{k\in\bb{N}}$ é um sistema projetivo de
feixes flácidos; demonstremos que satisfaz as condições da
parte~\emph{a)}.
\oldpage[0\textsubscript{III}]{69}
Com efeito, se $\sh{N}_{hk}$ denota o núcleo de $u_{hk}$ para $h\leq k$,
então a sucessão
\[
  0\longrightarrow\sh{N}_{hk}\longrightarrow\sh{F}_k
  \longrightarrow\sh{F}_h\longrightarrow0
\]
é exata por~{\rm(iii)}, e nossa asserção segue-se da exatidão
do funtor
$\sh{A}\longmapsto\sh{C}^i(X,\sh{A})$ (G, II, 4.3).

Ponhamos
\[
  \sh{G}^i=\varprojlim_k\sh{C}^i(X,\sh{F}_k).
\]
Pela parte~\emph{a)}, tem-se
$\HH^j(X,\sh{G}^i)=0$ para $j>0$ e $i\geq0$.
Demonstraremos que
$\sh{G}^\bullet=(\sh{G}^i)_{i\geq0}$ é uma resolução do feixe
$\sh{F}$. Como $\HH^j(X,\sh{G}^i)=0$ para $j>0$, a cohomologia
$\HH^\bullet(X,\sh{F})$ será então igual a
$\HH^\bullet(\Gamma(X,\sh{G}^\bullet))$ (G, II, 4.7.1).

Tomando o limite projetivo das sucessões exatas
\[
  0\longrightarrow\sh{F}_k
  \longrightarrow\sh{C}^0(X,\sh{F}_k)
  \longrightarrow\sh{C}^1(X,\sh{F}_k)
  \longrightarrow\cdots
\]
obtém-se manifestamente um complexo de feixes de grupos abelianos
\[
  0\longrightarrow\sh{F}
  \longrightarrow\sh{G}^0
  \longrightarrow\sh{G}^1
  \longrightarrow\cdots.
\]
Para demonstrar nossa asserção, resta estabelecer que
$\sh{H}^i(\sh{G}^\bullet)=0$ para $i>0$.
Este feixe é gerado pelo pré-feixe
\[
  U\longmapsto\HH^i\bigl(\Gamma(U,\sh{G}^\bullet)\bigr)
\]
(G, II, 4.1).
Ora, o complexo $\Gamma(U,\sh{G}^\bullet)$ é o limite projetivo
do sistema projetivo de complexos de grupos abelianos
\[
  \bigl(\Gamma(U,\sh{C}^\bullet(X,\sh{F}_k))\bigr)_{k\in\bb{N}}
\]
(0\textsubscript{I}, \EGAUnavailableHyperref{0.3.2.6}{3.2.6}).
Vimos na parte~\emph{a)} que, para todo $i\geq0$, as aplicações
\[
  \Gamma(U,\sh{C}^i(X,\sh{F}_k))
  \longrightarrow
  \Gamma(U,\sh{C}^i(X,\sh{F}_h))
  \qquad(h\leq k)
\]
são sobrejetivas. Por outro lado,
\[
  \HH^i(U,\sh{F}_k)
  =
  \HH^i\bigl(\Gamma(U,\sh{C}^\bullet(X,\sh{F}_k))\bigr),
\]
pois a resolução canônica $\sh{C}^\bullet(U,\sh{F}_k|U)$ é
induzida sobre $U$ por $\sh{C}^\bullet(X,\sh{F}_k)$.
Pela hipótese~{\rm(i)}, para todo $U\in\mathfrak{B}$ podemos aplicar
\sref{0.13.2.3} a este sistema projetivo de complexos, e portanto
\[
  \HH^i\bigl(\Gamma(U,\sh{G}^\bullet)\bigr)
  =
  \varprojlim_k\HH^i(U,\sh{F}_k)
  \qquad(i\geq0).
\]
A hipótese~{\rm(ii)} mostra então, por definição, que os feixes
$\sh{H}^i(\sh{G}^\bullet)$ se anulam para $i>0$.

Assim, para todo $i\geq0$,
\[
  \HH^i(X,\sh{F})
  =
  \HH^i\bigl(\Gamma(X,\sh{G}^\bullet)\bigr),
  \qquad
  \Gamma(X,\sh{G}^\bullet)
  =
  \varprojlim_k\Gamma(X,\sh{C}^\bullet(X,\sh{F}_k)).
\]
Acabamos de observar que todas as aplicações
\[
  \Gamma(X,\sh{C}^i(X,\sh{F}_k))
  \longrightarrow
  \Gamma(X,\sh{C}^i(X,\sh{F}_h))
  \qquad(h\leq k)
\]
são sobrejetivas; a conclusão segue-se, portanto, uma vez mais de
\sref{0.13.2.3}.
\end{proof}

\begin{env}[13.3.2]
\label{0.13.3.2}
\emph{Observações} (13.3.2).
\begin{enumerate}
  \item[{\rm(i)}]
  O enunciado de \sref{0.13.3.1} somente tem interesse para $i>0$,
  pois, para $i=0$, $h_i$ é sempre um isomorfismo sem hipótese
  alguma (0\textsubscript{I}, \EGAUnavailableHyperref{0.3.2.6}{3.2.6}).

  \item[{\rm(ii)}]
  As condições~{\rm(i)} e~{\rm(ii)} de \sref{0.13.3.1} são satisfeitas,
  em particular, se
  $\HH^i(U,\sh{F}_k)=0$ para todo $k$, todo $i>0$ e todo
  $U\in\mathfrak{B}$, e se, para $U\in\mathfrak{B}$, as aplicações
  \[
    \Gamma(U,\sh{F}_k)\longrightarrow\Gamma(U,\sh{F}_h)
  \]
  são sobrejetivas. Este é o caso mais frequente no qual se aplica
  \sref{0.13.3.1}.
\end{enumerate}
\end{env}

\subsection{Condição de Mittag-Leffler e objetos graduados associados a sistemas projetivos}
\label{subsection:0.13.4}
\label{0.13.4}

\begin{env}[13.4.1]
\label{0.13.4.1}
Seja
$\bb{A}=(A_k,u_{kh})_{k\in\bb{Z}}$ um sistema projetivo em uma
categoria abeliana $\cat{C}$.
Dizemos que está \emph{limitado inferiormente} se existe um $k_0$ tal que
$A_k=0$ para $k<k_0$.
Sobre cada $A_k$ definimos uma \emph{filtração}
$(\mathrm{F}^p(A_k))_{p\in\bb{Z}}$ por
\begin{equation}
\label{0.13.4.1.1}
\tag{13.4.1.1}
  \mathrm{F}^p(A_k)=
  \begin{cases}
    \Ker(A_k\longrightarrow A_{p-1}) & \text{para }p\leq k+1,\\
    0 & \text{para }p\geq k+1.
  \end{cases}
\end{equation}

\oldpage[0\textsubscript{III}]{70}
Por hipótese,
$\mathrm{F}^{k_0}(A_k)=A_k$ e
$\mathrm{F}^{k+1}(A_k)=0$; em outras palavras, a filtração considerada
é finita (\sref{0.11.1.3}).
Os objetos graduados associados a esta filtração são, portanto,
\[
  \mathrm{gr}^p(A_k)
  =
  \Ker(A_k\longrightarrow A_{p-1})
  \big/
  \Ker(A_k\longrightarrow A_p).
\]
Por conseguinte, $\mathrm{gr}^p(A_k)$ é isomorfo à imagem por
$A_k\longrightarrow A_p$ de
$\Ker(A_k\longrightarrow A_{p-1})$.
Pela transitividade dos morfismos que definem um sistema projetivo,
tem-se, portanto,
\begin{equation}
\label{0.13.4.1.2}
\tag{13.4.1.2}
  \mathrm{gr}^p(A_k)
  =
  \Ker(A_p\longrightarrow A_{p-1})
  \cap
  \Im(A_k\longrightarrow A_p).
\end{equation}
Mas \sref{0.13.4.1.1} dá
$\Ker(A_p\longrightarrow A_{p-1})=\mathrm{gr}^p(A_p)$, e portanto
\begin{equation}
\label{0.13.4.1.3}
\tag{13.4.1.3}
  \mathrm{gr}^p(A_k)
  =
  \mathrm{gr}^p(A_p)
  \cap
  \Im(A_k\longrightarrow A_p).
\end{equation}

As definições precedentes mostram também que, para $k\leq h$,
\[
  u_{kh}\bigl(\mathrm{F}^p(A_h)\bigr)
  \subset
  \mathrm{F}^p(A_k),
\]
e, por conseguinte, para todo $p\in\bb{Z}$, os
$\mathrm{gr}^p(u_{kh})$ definem um sistema projetivo
$(\mathrm{gr}^p(A_k))_{k\in\bb{Z}}$.
\end{env}

\begin{env}[13.4.2]
\label{0.13.4.2}
Dizemos que o sistema projetivo $\bb{A}$ é
\emph{essencialmente constante} se os morfismos
$A_{k+1}\longrightarrow A_k$ são isomorfismos para todo $k$
suficientemente grande.
Dizemos que $\bb{A}$ é \emph{estrito} se os morfismos
$A_i\longrightarrow A_j$ $(j\leq i)$ são epimorfismos.

Quando $\bb{A}$ é estrito, segue-se de \sref{0.13.4.1.3} que, para
$p\leq k\leq h$, o morfismo canônico
\[
  \mathrm{gr}^p(A_h)\longrightarrow\mathrm{gr}^p(A_k)
\]
é um isomorfismo; em outras palavras, o sistema projetivo
$(\mathrm{gr}^p(A_k))_{k\in\bb{Z}}$ é essencialmente constante.
A sucessão de objetos $\mathrm{gr}^p(A_p)$, identificada para todo
$p$ com $\varprojlim_k\mathrm{gr}^p(A_k)$, denota-se então
$\mathrm{gr}^\bullet(\bb{A})$ e chama-se o
\emph{objeto graduado associado ao sistema projetivo estrito}
$\bb{A}=(A_k)$.

Suponhamos agora que o sistema projetivo $\bb{A}$, limitado
inferiormente, satisfaz (ML).
Por \sref{0.13.1.2}, o sistema projetivo
$\bb{A}'=(A'_k)$ de objetos de ``imagem universal'' é estrito; também está
limitado inferiormente.
O objeto graduado $\mathrm{gr}^\bullet(\bb{A}')$ associado a
$\bb{A}'$ chama-se igualmente o objeto graduado
\emph{associado a} $\bb{A}$ e é denotado
$\mathrm{gr}^\bullet(\bb{A})$.
\end{env}

\begin{proposition}[13.4.3]
\label{0.13.4.3}
Seja $\bb{A}=(A_k)_{k\in\bb{Z}}$ um sistema projetivo limitado
inferiormente em uma categoria abeliana $\cat{C}$.
As condições seguintes são equivalentes:
\begin{enumerate}
  \item[{\rm a)}] $\bb{A}$ satisfaz (ML).
  \item[{\rm b)}] Para todo $p\in\bb{Z}$, o sistema projetivo
  $(\mathrm{gr}^p(A_k))_{k\in\bb{Z}}$ é essencialmente constante.
\end{enumerate}
Além disso, quando essas condições são satisfeitas, existe, para todo
$p\in\bb{Z}$, um isomorfismo canônico
\begin{equation}
\label{0.13.4.3.1}
\tag{13.4.3.1}
  \mathrm{gr}^p(\bb{A})
  \isoto
  \varprojlim_k\mathrm{gr}^p(A_k).
\end{equation}
\end{proposition}

\begin{proof}
Segue-se imediatamente de \sref{0.13.4.1.2} que \emph{a)} implica
\emph{b)}. A mesma fórmula, aplicada ao sistema projetivo
$\bb{A}'$ com a notação de \sref{0.13.4.2}, da
\sref{0.13.4.3.1} por definição.

Para $k\leq h$, ponhamos
\[
  A_{kh}=\Im(A_h\longrightarrow A_k).
\]
Se $k\leq h\leq j$, então
$A_{kj}\subset A_{kh}\subset A_k$.
Dotemos $A_{kh}$ da filtração induzida por
$(\mathrm{F}^p(A_k))$.
A transitividade dos morfismos que definem $\bb{A}$ mostra
imediatamente que esta é também a filtração quociente de
$(\mathrm{F}^p(A_h))$; por conseguinte,
\begin{equation}
\label{0.13.4.3.2}
\tag{13.4.3.2}
  \mathrm{gr}^p(A_{kh})
  =
  \Im\bigl(
    \mathrm{gr}^p(A_h)\longrightarrow\mathrm{gr}^p(A_k)
  \bigr).
\end{equation}

\oldpage[0\textsubscript{III}]{71}
Suponhamos agora que \emph{b)} seja satisfeita.
Para todo $p\in\bb{Z}$ e todo $k\geq p$, existe um inteiro
$\mathrm{L}(p,k)$ tal que o membro direito de
\sref{0.13.4.3.2} é constante para $h\geq\mathrm{L}(p,k)$.
Como $\mathrm{gr}^p(A_k)=0$ para $p<k_0$, para $k$ fixo somente
um número finito de inteiros $\mathrm{L}(p,k)$ pode ser não nulo quando
$p$ percorre os inteiros $\leq k$.
Seja $\mathrm{L}(k)=m$ o maior destes inteiros.
Para todo $h\geq m$, tem-se
$A_{kh}\subset A_{km}$, e, pela definição de $m$, a injeção
canônica
$A_{kh}\longrightarrow A_{km}$ induz um isomorfismo
\[
  \mathrm{gr}^\bullet(A_{kh})
  \simto
  \mathrm{gr}^\bullet(A_{km}).
\]
Como as filtrações são finitas, esta injeção é ela própria bijetiva
(Bourbaki, \emph{Alg.\ comm.}, cap.~III, \textsection2,
n\textsuperscript{o}~8, teorema~1).
Isto demonstra que $\bb{A}$ satisfaz (ML).
\end{proof}

\begin{env}[13.4.4]
\label{0.13.4.4}
Suponhamos que o limite projetivo
$A=\varprojlim_k A_k$ existe em $\cat{C}$.
Nas definições de \sref{0.13.4.1}, se pode então substituir $A_k$
por $A$, e a filtração assim definida sobre $A$ continua satisfazendo
\begin{equation}
\label{0.13.4.4.1}
\tag{13.4.4.1}
  \mathrm{gr}^p(A)
  =
  \mathrm{gr}^p(A_p)
  \cap
  \Im(A\longrightarrow A_p).
\end{equation}
\end{env}

\begin{corollary}[13.4.5]
\label{0.13.4.5}
Suponhamos que $\cat{C}$ é a categoria de grupos abelianos.
Se o sistema projetivo $\bb{A}$ satisfaz (ML) e
$A=\varprojlim_k A_k$, então para todo $p\in\bb{Z}$ existe um
isomorfismo canônico
\begin{equation}
\label{0.13.4.5.1}
\tag{13.4.5.1}
  \mathrm{gr}^p(A)
  \isoto
  \varprojlim_k\mathrm{gr}^p(A_k).
\end{equation}
\end{corollary}

\begin{proof}
Com efeito,
\[
  \Im(A_k\longrightarrow A_p)
  =
  \Im(A\longrightarrow A_p)
\]
para todo $k$ suficientemente grande
(Bourbaki, \emph{Top.\ g\'{e}n.}, cap.~II,
3\textsuperscript{e}~\'{e}d., \textsection3,
n\textsuperscript{o}~5, teorema~1), e a conclusão decorre de
\sref{0.13.4.1.3} e \sref{0.13.4.4.1}.
\end{proof}

\subsection{Limites projetivos de sucessões espectrais de complexos filtrados}
\label{subsection:0.13.5}

\begin{env}[13.5.1]
\label{0.13.5.1}
Seja $\cat{C}$ uma categoria abeliana e seja $X^\bullet$ um complexo
de objetos de $\cat{C}$ dotado de uma filtração
$(\mathrm{F}^p(X^\bullet))_{p\in\bb{Z}}$ tal que
$\mathrm{F}^{p_0}(X^\bullet)=X^\bullet$ para algum índice $p_0$.
Para todo $k\in\bb{Z}$, consideremos o complexo
\[
  X^\bullet_k=X^\bullet/\mathrm{F}^{k+1}(X^\bullet);
\]
está dotado canonicamente da filtração constituída por
\[
  \mathrm{F}^p(X^\bullet_k)
  =
  \mathrm{F}^p(X^\bullet)/\mathrm{F}^{k+1}(X^\bullet)
  \quad\text{para }p\leq k,
\]
e $\mathrm{F}^p(X^\bullet_k)=0$ para $p\geq k+1$.
Além disso, existem morfismos canônicos
$X^\bullet_{k+1}\longrightarrow X^\bullet_k$, que fazem de
$\bb{X}^\bullet=(X^\bullet_k)_{k\in\bb{Z}}$ um \emph{sistema projetivo
de complexos filtrados de objetos de $\cat{C}$}.
Observemos que este sistema projetivo é \emph{estrito} e que
$X^\bullet_k=0$ para $k<p_0$.
\end{env}

\begin{env}[13.5.2]
\label{0.13.5.2}
Mais geralmente, consideremos um sistema projetivo estrito limitado
inferiormente
$\bb{X}^\bullet=(X^\bullet_k)_{k\in\bb{Z}}$ de complexos de objetos de
$\cat{C}$, e dotemos cada $X^\bullet_k$ da filtração definida em
\sref{0.13.4.1} (aplicada na categoria abeliana de complexos limitados
inferiormente de objetos de $\cat{C}$).
Os morfismos
$X^\bullet_k\longrightarrow X^\bullet_p$ ($p\leq k$) tornam-se então
em morfismos de complexos \emph{filtrados} com filtrações finitas.
A funtorialidade das sucessões espectrais de complexos filtrados
(\sref{0.11.2.3}) mostra que os morfismos de transição de
$\bb{X}^\bullet$ induzem morfismos que fazem de
\[
  \mathrm{E}(\bb{X}^\bullet)
  =
  \bigl(\mathrm{E}(X^\bullet_k)\bigr)_{k\in\bb{Z}}
\]
um \emph{sistema projetivo de sucessões espectrais}.
\end{env}

\begin{lemma}[13.5.3]
\label{0.13.5.3}
Suponhamos que o sistema projetivo
$\bb{X}^\bullet=(X^\bullet_k)_{k\in\bb{Z}}$ de complexos filtrados se
obtém como em \sref{0.13.5.2}. Então:
\begin{enumerate}
  \item[{\rm a)}] Para $r\geq p-p_0$, tem-se
  \[
    \mathrm{B}_r^{pq}(X^\bullet_k)
    =
    \mathrm{B}_\infty^{pq}(X^\bullet_k)
  \]
  para todo $k\in\bb{Z}$.
  \item[{\rm b)}] Para $k+1\leq p+r$, tem-se
  \[
    \mathrm{Z}_r^{pq}(X^\bullet_k)
    =
    \mathrm{Z}_\infty^{pq}(X^\bullet_k).
  \]
  \item[{\rm c)}] Para $k+1\geq p+r$, os morfismos
  \[
    \mathrm{Z}_r^{pq}(X^\bullet_h)
    \longrightarrow
    \mathrm{Z}_r^{pq}(X^\bullet_k)
    \quad\text{e}\quad
    \mathrm{B}_r^{pq}(X^\bullet_h)
    \longrightarrow
    \mathrm{B}_r^{pq}(X^\bullet_k)
  \]
  são isomorfismos para todo $h\geq k$.
\end{enumerate}
\end{lemma}

\oldpage[0\textsubscript{III}]{72}
Estas três propriedades seguem-se imediatamente das definições de
\sref{0.11.2.2}, levando em conta que
$\mathrm{F}^{p-r+1}(X^\bullet_k)=X^\bullet_k$ quando $p-r<p_0$.

\begin{env}[13.5.4]
\label{0.13.5.4}
Suponhamos que sejam satisfeitas as hipóteses de \sref{0.13.5.3}.
Então, para $p$, $q$ e $r$ fixos (com $r$ finito), os sistemas
projetivos
\[
  \bigl(\mathrm{Z}_r^{pq}(X^\bullet_k)\bigr)_{k\in\bb{Z}},
  \qquad
  \bigl(\mathrm{B}_r^{pq}(X^\bullet_k)\bigr)_{k\in\bb{Z}},
  \qquad
  \bigl(\mathrm{E}_r^{pq}(X^\bullet_k)\bigr)_{k\in\bb{Z}}
\]
são \emph{essencialmente constantes}; seus respectivos limites projetivos
serão denotados por
$\mathrm{Z}_r^{pq}(\bb{X}^\bullet)$,
$\mathrm{B}_r^{pq}(\bb{X}^\bullet)$, e
\[
  \mathrm{E}_r^{pq}(\bb{X}^\bullet)
  =
  \mathrm{Z}_r^{pq}(\bb{X}^\bullet)/
  \mathrm{B}_r^{pq}(\bb{X}^\bullet).
\]
Os objetos $\mathrm{Z}_r^{pq}(\bb{X}^\bullet)$ e
$\mathrm{B}_r^{pq}(\bb{X}^\bullet)$ identificam-se canonicamente com
subobjetos de $\mathrm{E}_1^{pq}(\bb{X}^\bullet)$.
A definição dos $d_r^{pq}$ (M, XV, 1) mostra que estes
morfismos, tomados com respeito aos $X^\bullet_k$, são também essencialmente
constantes, e definem, portanto, morfismos
\begin{equation}
\label{0.13.5.4.1}
\tag{13.5.4.1}
  d_r^{pq}:
  \mathrm{E}_r^{pq}(\bb{X}^\bullet)
  \longrightarrow
  \mathrm{E}_r^{p+r,q-r+1}(\bb{X}^\bullet)
\end{equation}
tal que
$d_r^{p+r,q-r+1}\circ d_r^{pq}=0$.
Além disso, existem isomorfismos canônicos de
$\Ker(d_r^{pq})$ sobre
\[
  \mathrm{Z}_{r+1}^{pq}(\bb{X}^\bullet)/
  \mathrm{B}_r^{pq}(\bb{X}^\bullet)
\]
e de $\Im(d_r^{pq})$ sobre
\[
  \mathrm{B}_{r+1}^{p+r,q-r+1}(\bb{X}^\bullet)/
  \mathrm{B}_r^{p+r,q-r+1}(\bb{X}^\bullet).
\]
\end{env}

\begin{lemma}[13.5.5]
\label{0.13.5.5}
Sob as hipóteses de \sref{0.13.5.3}, para
$s\geq r>p-p_0$ existe um monomorfismo canônico
\begin{equation}
\label{0.13.5.5.1}
\tag{13.5.5.1}
  i:
  \mathrm{E}_s^{pq}(\bb{X}^\bullet)
  \longrightarrow
  \mathrm{E}_r^{pq}(\bb{X}^\bullet)
\end{equation}
e um isomorfismo canônico
\begin{equation}
\label{0.13.5.5.2}
\tag{13.5.5.2}
  j_r:
  \mathrm{E}_r^{pq}(\bb{X}^\bullet)
  \simto
  \mathrm{E}_\infty^{pq}(X^\bullet_{p+r-1})
\end{equation}
tal que o diagrama
\begin{equation}
\label{0.13.5.5.3}
\tag{13.5.5.3}
  \xymatrix{
    \mathrm{E}_s^{pq}(\bb{X}^\bullet)
      \ar[r]^{j_s}\ar[d]_i
    &
    \mathrm{E}_\infty^{pq}(X^\bullet_{p+s-1})
      \ar[d]
    \\
    \mathrm{E}_r^{pq}(\bb{X}^\bullet)
      \ar[r]^{j_r}
    &
    \mathrm{E}_\infty^{pq}(X^\bullet_{p+r-1})
  }
\end{equation}
é comutativo, onde a seta vertical direita está induzida pelo
morfismo
$X^\bullet_{p+s-1}\longrightarrow X^\bullet_{p+r-1}$.
\end{lemma}

\begin{proof}
A existência de $i$ segue-se de
\[
  \mathrm{B}_r^{pq}(X^\bullet_k)
  =
  \mathrm{B}_\infty^{pq}(X^\bullet_k)
  \qquad (r>p-p_0)
\]
(\sref{0.13.5.3},~a)).
Para $k+1\leq p+r$, tem-se
\[
  \mathrm{Z}_r^{pq}(X^\bullet_k)
  =
  \mathrm{Z}_\infty^{pq}(X^\bullet_k)
\]
(\sref{0.13.5.3},~b)); em particular,
\[
  \mathrm{Z}_r^{pq}(X^\bullet_{p+r-1})
  =
  \mathrm{Z}_\infty^{pq}(X^\bullet_{p+r-1}).
\]
Por outro lado,
$\mathrm{Z}_r^{pq}(X^\bullet_{p+r-1})$ e
$\mathrm{B}_r^{pq}(X^\bullet_{p+r-1})$ identificam-se canonicamente com
$\mathrm{Z}_r^{pq}(\bb{X}^\bullet)$ e
$\mathrm{B}_r^{pq}(\bb{X}^\bullet)$ por
\sref{0.13.5.3},~c).
Isto demonstra a existência de $j_r$ e a comutatividade de
\sref{0.13.5.5.3}.
\end{proof}

\begin{corollary}[13.5.6]
\label{0.13.5.6}
Sob as hipóteses de \sref{0.13.5.3}, se um dos limites
projetivos
\[
  \varprojlim_r\mathrm{E}_r^{pq}(\bb{X}^\bullet),
  \qquad
  \varprojlim_k\mathrm{E}_\infty^{pq}(X^\bullet_k)
\]
existe, então também existe o outro, e há um isomorfismo canônico
\begin{equation}
\label{0.13.5.6.1}
\tag{13.5.6.1}
  j_\infty:
  \varprojlim_r\mathrm{E}_r^{pq}(\bb{X}^\bullet)
  \simto
  \varprojlim_k\mathrm{E}_\infty^{pq}(X^\bullet_k).
\end{equation}
Além disso, o sistema projetivo
$(\mathrm{E}_r^{pq}(\bb{X}^\bullet))_{r\in\bb{Z}}$ é essencialmente
constante (\sref{0.13.4.2}) se e somente se o sistema projetivo
$(\mathrm{E}_\infty^{pq}(X^\bullet_k))_{k\in\bb{Z}}$ é essencialmente
constante.
\end{corollary}

\begin{env}[13.5.7]
\label{0.13.5.7}
Denotamos por
$\mathrm{B}_\infty^{pq}(\bb{X}^\bullet)$ e
$\mathrm{Z}_\infty^{pq}(\bb{X}^\bullet)$ os subobjetos de
$\mathrm{E}_1^{pq}(\bb{X}^\bullet)$ iguais respectivamente a
$\mathrm{B}_r^{pq}(\bb{X}^\bullet)$ para $r>p-p_0$ e a
$\inf_r\mathrm{Z}_r^{pq}(\bb{X}^\bullet)$, quando este último existe.
Assim,
\[
  \varprojlim_r\mathrm{E}_r^{pq}(\bb{X}^\bullet)
\]
\oldpage[0\textsubscript{III}]{73}
identifica-se canonicamente com
\[
  \mathrm{E}_\infty^{pq}(\bb{X}^\bullet)
  =
  \mathrm{Z}_\infty^{pq}(\bb{X}^\bullet)/
  \mathrm{B}_\infty^{pq}(\bb{X}^\bullet).
\]
Observemos que os objetos
$\mathrm{Z}_r^{pq}(\bb{X}^\bullet)$,
$\mathrm{B}_r^{pq}(\bb{X}^\bullet)$,
$\mathrm{E}_r^{pq}(\bb{X}^\bullet)$
($1\leq r\leq+\infty$), e os morfismos $d_r^{pq}$ dependem
\emph{funtorialmente} do sistema projetivo $\bb{X}^\bullet$, com as
restrições de \sref{0.13.5.5}, e os morfismos definidos em
\sref{0.13.5.5} e \sref{0.13.5.6} são funtoriais.
\end{env}

\subsection{Sucessão espectral de um funtor relativa a um objeto dotado de uma filtração finita}
\label{subsection:0.13.6}

\begin{env}[13.6.1]
\label{0.13.6.1}
Sejam $\cat{C}$ e $\cat{C}'$ duas categorias abelianas, e seja
$\mathrm{T}:\cat{C}\longrightarrow\cat{C}'$ um funtor covariante
aditivo.
Suponhamos que todo objeto de $\cat{C}$ é isomorfo a um subobjeto de
um objeto injetivo, de modo que existam os funtores derivados à direita
$\RR^p\mathrm{T}$ ($p\geq0$).
\end{env}

\begin{lemma}[13.6.2]
\label{0.13.6.2}
Seja $A$ um objeto de $\cat{C}$ dotado de uma filtração finita
$(\mathrm{F}^i(A))_{i\in\bb{Z}}$.
Existe uma resolução injetiva
$X^\bullet=(X^j)_{j\geq0}$ de $A$, dotada de uma filtração finita
$(\mathrm{F}^i(X^\bullet))_{i\in\bb{Z}}$, tal que
\[
  \mathrm{F}^i(A)=A
  \quad\text{(respectivamente, }\mathrm{F}^i(A)=0\text{)}
\]
implica
\[
  \mathrm{F}^i(X^\bullet)=X^\bullet
  \quad\text{(respectivamente, }\mathrm{F}^i(X^\bullet)=0\text{)},
\]
e tal que, para todo $i\in\bb{Z}$,
$\mathrm{F}^i(X^\bullet)$ é uma resolução injetiva de
$\mathrm{F}^i(A)$.
\end{lemma}

\begin{proof}
Seja $p$ (respectivamente, $q>p$) o maior índice tal que
$\mathrm{F}^p(A)=A$ (respectivamente, o menor índice tal que
$\mathrm{F}^q(A)=0$).
Raciocinamos por indução sobre $q-p$, sendo imediato o lema quando
$q-p=1$.
Suponhamos construída uma resolução injetiva ${X''}^\bullet$ de
$A/\mathrm{F}^{q-1}(A)$ que tenha as propriedades requeridas.
Consideremos a sucessão exata
\[
  0\longrightarrow\mathrm{F}^{q-1}(A)
  \longrightarrow A
  \longrightarrow A/\mathrm{F}^{q-1}(A)
  \longrightarrow0,
\]
tomemos uma resolução injetiva ${X'}^\bullet$ de
$\mathrm{F}^{q-1}(A)$, e escolhamos depois uma resolução injetiva
$X^\bullet$ de $A$ de modo que se obtenha uma sucessão exata
\[
  0\longrightarrow {X'}^\bullet
  \longrightarrow X^\bullet
  \longrightarrow {X''}^\bullet
  \longrightarrow0
\]
compatível com a sucessão precedente (M, V, 2.2).
Está claro que $X^\bullet$ tem as propriedades requeridas.
\end{proof}

\begin{corollary}[13.6.3]
\label{0.13.6.3}
Seja $B$ um segundo objeto de $\cat{C}$, dotado de uma filtração
finita $(\mathrm{F}^i(B))_{i\in\bb{Z}}$, seja $s$ um inteiro,
e seja $u:A\longrightarrow B$ um morfismo tal que
\[
  u(\mathrm{F}^i(A))\subset\mathrm{F}^{i+s}(B)
\]
para todo $i\in\bb{Z}$.
Se $Y^\bullet=(Y^j)_{j\geq0}$ é uma resolução injetiva de $B$
dotada de uma filtração
$(\mathrm{F}^i(Y^\bullet))_{i\in\bb{Z}}$ que tenha as propriedades
enunciadas em \sref{0.13.6.2}, então existe um morfismo
$v:X^\bullet\longrightarrow Y^\bullet$, compatível com $u$, tal
que
\[
  v(\mathrm{F}^i(X^\bullet))
  \subset
  \mathrm{F}^{i+s}(Y^\bullet)
\]
para todo $i\in\bb{Z}$.
Além disso, quaisquer dois desses morfismos $v$ e $v'$ são homotópicos.
\end{corollary}

Isto segue-se imediatamente, por indução sobre $q-p$, da construção
precedente e de (M, V, 2.3).

\begin{env}[13.6.4]
\label{0.13.6.4}
Sob as hipóteses de \sref{0.13.6.2}, consideremos o complexo
$\mathrm{T}(X^\bullet)$ em $\cat{C}'$.
Está filtrado pelos complexos
$\mathrm{T}(\mathrm{F}^i(X^\bullet))$, pois
$\mathrm{F}^i(X^\bullet)$ é um sumando direto de $X^\bullet$.
Segue-se de \sref{0.13.6.3} que a \emph{sucessão espectral} de
este complexo filtrado depende, salvo isomorfismo, unicamente do
objeto filtrado $A$.
Seu termo limite é a cohomologia $\RR^\bullet\mathrm{T}(A)$, dotada
da filtração
\begin{equation}
\label{0.13.6.4.1}
\tag{13.6.4.1}
  \begin{split}
  \mathrm{F}^p(\RR^n\mathrm{T}(A))
  &=
  \Im\bigl(
    \RR^n\mathrm{T}(\mathrm{F}^p(A))
    \longrightarrow
    \RR^n\mathrm{T}(A)
  \bigr)
  \\
  &=
  \Ker\bigl(
    \RR^n\mathrm{T}(A)
    \longrightarrow
    \RR^n\mathrm{T}(A/\mathrm{F}^p(A))
  \bigr),
  \end{split}
\end{equation}
por \sref{0.11.2.2}, e seu termo $\mathrm{E}_1$ é
\begin{equation}
\label{0.13.6.4.2}
\tag{13.6.4.2}
  \mathrm{E}_1^{pq}
  =
  \RR^{p+q}\mathrm{T}(\mathrm{gr}^p(A)),
\end{equation}
onde, como de costume,
$\mathrm{gr}^p(A)=\mathrm{F}^p(A)/\mathrm{F}^{p+1}(A)$.
Deduz-se de \sref{0.11.2.2} que a filtração do termo limite
é finita e que, para $p$ e $q$ fixos, as sucessões
\oldpage[0\textsubscript{III}]{74}
\[
  \mathrm{B}_r^{pq}(A)
  =
  \mathrm{B}_r^{pq}(\mathrm{T}(X^\bullet))
  \quad\text{e}\quad
  \mathrm{Z}_r^{pq}(A)
  =
  \mathrm{Z}_r^{pq}(\mathrm{T}(X^\bullet))
\]
são estacionárias.
A sucessão espectral precedente é, portanto, \emph{birregular}
(\sref{0.11.1.3}).
Denotamo-la por
$\mathrm{E}(A)=(\mathrm{E}_r^{pq}(A))$ e chamamo-la de
\emph{sucessão espectral do funtor $\mathrm{T}$ relativa ao
objeto filtrado $A$}.
\end{env}

\begin{env}[13.6.5]
\label{0.13.6.5}
Suponhamos agora que sejam satisfeitas as hipóteses de \sref{0.13.6.3}, e
conservemos sua notação.
Como $\mathrm{F}^i(X^\bullet)$ (respectivamente,
$\mathrm{F}^i(Y^\bullet)$) é um sumando direto de $X^\bullet$
(respectivamente, $Y^\bullet$), tem-se
\[
  (\mathrm{T}v)\bigl(\mathrm{T}(\mathrm{F}^i(X^\bullet))\bigr)
  \subset
  \mathrm{T}(\mathrm{F}^{i+s}(Y^\bullet))
\]
para todo $i\in\bb{Z}$.
As definições de \sref{0.11.2.2} mostram então que, para
$1\leq r\leq+\infty$, $\mathrm{T}v$ define morfismos
\[
  \mathrm{B}_r^{pq}(\mathrm{T}(X^\bullet))
  \longrightarrow
  \mathrm{B}_r^{p+s,q-s}(\mathrm{T}(Y^\bullet))
\]
e
\[
  \mathrm{Z}_r^{pq}(\mathrm{T}(X^\bullet))
  \longrightarrow
  \mathrm{Z}_r^{p+s,q-s}(\mathrm{T}(Y^\bullet)),
\]
e, portanto, um morfismo
\[
  w_r:
  \mathrm{E}_r^{pq}(A)
  \longrightarrow
  \mathrm{E}_r^{p+s,q-s}(B).
\]
Analogamente, sobre os termos limite têm-se morfismos
\[
  u_n:
  \RR^n\mathrm{T}(A)
  \longrightarrow
  \RR^n\mathrm{T}(B)
\]
tais que
\[
  u_n\bigl(\mathrm{F}^p(\RR^n\mathrm{T}(A))\bigr)
  \subset
  \mathrm{F}^{p+s}(\RR^n\mathrm{T}(B)).
\]
A definição dos $d_r^{pq}$ (M, XV, 1) mostra também que os
diagramas
\[
  \xymatrix{
    \mathrm{E}_r^{pq}(A)
      \ar[r]^{d_r^{pq}}\ar[d]_{w_r}
    &
    \mathrm{E}_r^{p+r,q-r+1}(A)
      \ar[d]^{w_r}
    \\
    \mathrm{E}_r^{p+s,q-s}(B)
      \ar[r]_{d_r^{p+s,q-s}}
    &
    \mathrm{E}_r^{p+r+s,q-r-s+1}(B)
  }
\]
são comutativos.
Segue-se que existe um diagrama comutativo análogo para os
isomorfismos $\alpha_r^{pq}$, cuja explicitação deixamos ao leitor.
Finalmente (\emph{loc.\ cit.}), têm-se também diagramas comutativos
sobre os termos limite:
\[
  \xymatrix{
    \mathrm{E}_\infty^{pq}(A)
      \ar[r]^{\beta^{pq}}\ar[d]_{w_\infty}
    &
    \mathrm{gr}^p(\RR^{p+q}\mathrm{T}(A))
      \ar[d]^{u_{p+q}}
    \\
    \mathrm{E}_\infty^{p+s,q-s}(B)
      \ar[r]_{\beta^{p+s,q-s}}
    &
    \mathrm{gr}^{p+s}(\RR^{p+q}\mathrm{T}(B)).
  }
\]
\end{env}

\begin{env}[13.6.6]
\label{0.13.6.6}
Suponhamos, em particular, que existe um anel $S$, dotado de uma
\emph{filtração} $(\mathrm{F}^i(S))_{i\in\bb{Z}}$, e um homomorfismo
de anéis
\begin{equation}
\label{0.13.6.6.1}
\tag{13.6.6.1}
  h:S\longrightarrow\Hom_{\cat{C}}(A,A)
\end{equation}
\oldpage[0\textsubscript{III}]{75}
tal que, para todo $t\in\mathrm{F}^i(S)$,
\[
  h_t(\mathrm{F}^j(A))
  \subset
  \mathrm{F}^{i+j}(A)
\]
para todo par $i,j$.
Para abreviar, dizemos então que $A$ está dotado da estrutura de um
\emph{$S$--$\cat{C}$-módulo filtrado} sobre o anel filtrado $S$.

Ao passar aos objetos graduados associados, todo $h_t$, para
$t\in\mathrm{F}^i(S)$, define um endomorfismo graduado
$\overline{h}_t$ de $\mathrm{gr}^\bullet(A)$, homogêneo de grau
$i$.
Além disso, este endomorfismo depende unicamente da classe de $t$ em
$\mathrm{gr}^i(S)$, e obtém-se assim um homomorfismo de
\emph{anéis graduados}
\[
  \overline{h}:
  \mathrm{gr}^\bullet(S)
  \longrightarrow
  \Hom_{\cat{C}}^\bullet
  \bigl(\mathrm{gr}^\bullet(A),\mathrm{gr}^\bullet(A)\bigr),
\]
onde o membro direito é o anel de endomorfismos graduados de
$\mathrm{gr}^\bullet(A)$.
Dizemos que $\mathrm{gr}^\bullet(A)$ está dotado da estrutura
de um \emph{$\mathrm{gr}^\bullet(S)$--$\cat{C}$-módulo graduado}.

Segue-se de \sref{0.13.6.5} que, para
$1\leq r\leq+\infty$, todo
$\overline{t}\in\mathrm{gr}^i(S)$ define canonicamente, sobre os
objetos bigraduados
\[
  \bigl(\mathrm{B}_r^{pq}(A)\bigr)_{(p,q)\in\bb{Z}\times\bb{Z}},
  \qquad
  \bigl(\mathrm{Z}_r^{pq}(A)\bigr)_{(p,q)\in\bb{Z}\times\bb{Z}},
\]
e
\[
  \mathrm{E}_r(A)
  =
  \bigl(\mathrm{E}_r^{pq}(A)\bigr)_{(p,q)\in\bb{Z}\times\bb{Z}},
\]
endomorfismos bigraduados de grau $(i,-i)$.
Sobre $\mathrm{E}_r(A)$, para $r$ finito, este endomorfismo comuta com
o endomorfismo bigraduado definido pelos $d_r^{pq}$.
Estes endomorfismos satisfazem as condições usuais de associatividade e
distributividade com respeito à adição em
$\mathrm{gr}^\bullet(S)$ e nos objetos bigraduados em questão.
Para abreviar, dizemos que estes últimos são
\emph{$\mathrm{gr}^\bullet(S)$--$\cat{C}'$-módulos bigraduados};
é imediato que os $\alpha_r^{pq}$ definem um isomorfismo para
este tipo de estrutura.

Para todo inteiro $n$, seja $\mathrm{B}_r^{(n)}(A)$
(respectivamente, $\mathrm{Z}_r^{(n)}(A)$ e
$\mathrm{E}_r^{(n)}(A)$) o subobjeto graduado de
$\mathrm{B}_r^{\bullet,\bullet}(A)$
(respectivamente, $\mathrm{Z}_r^{\bullet,\bullet}(A)$ e
$\mathrm{E}_r^{\bullet,\bullet}(A)$) constituído pelos termos para
os quais $p+q=n$ ($1\leq r\leq+\infty$).
É imediato que estes são
\emph{$\mathrm{gr}^\bullet(S)$--$\cat{C}'$-módulos graduados}.

Finalmente, todo $\overline{t}\in\mathrm{gr}^i(S)$ define, para todo
$n$, um endomorfismo graduado de grau $i$ sobre o objeto graduado
$\mathrm{gr}^\bullet(\RR^n\mathrm{T}(A))$, que adquire assim a
estrutura de um
$\mathrm{gr}^\bullet(S)$--$\cat{C}'$-módulo graduado.
Para $p+q=n$, os $\beta^{pq}$ definem um isomorfismo
\[
  \mathrm{E}^{(n)}(A)
  \isoto
  \mathrm{gr}^\bullet(\RR^n\mathrm{T}(A))
\]
para este tipo de estrutura.

Observemos que, quando $\cat{C}'$ é a categoria de grupos abelianos, as
estruturas de $S$--$\cat{C}'$-módulo (respectivamente, de
$\mathrm{gr}^\bullet(S)$--$\cat{C}'$-módulo graduado ou bigraduado) são
simplesmente as estruturas usuais de $S$-módulo (respectivamente, de
$\mathrm{gr}^\bullet(S)$-módulo graduado ou bigraduado).
\end{env}

\subsection{Funtores derivados de um limite projetivo de argumentos}
\label{subsection:0.13.7}
\label{0.13.7}

\begin{env}[13.7.1]
\label{0.13.7.1}
Sejam $\cat{C}$ e $\cat{C}'$ duas categorias abelianas, e suponhamos
que todo objeto de $\cat{C}$ é isomorfo a um subobjeto de um
objeto injetivo.
Seja
\[
  \mathrm{T}:\cat{C}\longrightarrow\cat{C}'
\]
um funtor covariante aditivo.
Consideremos um sistema projetivo estrito
$\bb{A}=(A_k)_{k\in\bb{Z}}$ em $\cat{C}$, limitado inferiormente; mais
precisamente, suponhamos que $A_k=0$ para $k<k_0$.
A este sistema associamos canonicamente, sobre cada $A_k$, a
filtração $(\mathrm{F}^p(A_k))_{p\in\bb{Z}}$ dada por
\sref{0.13.4.1.1}.
Como o sistema projetivo é estrito, os morfismos canônicos
\begin{equation}
\label{0.13.7.1.1}
\tag{13.7.1.1}
  \mathrm{F}^i(A_h)/\mathrm{F}^j(A_h)
  \longrightarrow
  \mathrm{F}^i(A_k)/\mathrm{F}^j(A_k)
  \qquad(h\geq k)
\end{equation}
são isomorfismos para $i\leq j\leq k+1$.
Lembremos também que
$\mathrm{F}^{k_0}(A_k)=A_k$ e
$\mathrm{F}^{k+1}(A_k)=0$ para todo $k$.
\end{env}

\begin{env}[13.7.2]
\label{0.13.7.2}
Para cada $k$, construamos uma resolução injetiva
$\mathrm{X}_k^\bullet=(\mathrm{X}_k^j)_{j\geq0}$ de $A_k$ que tenha as
propriedades de \sref{0.13.6.2}.
Os morfismos canônicos $A_{k+1}\longrightarrow A_k$ nos permitem, por
\sref{0.13.6.3}, definir para cada $k$ um morfismo de complexos
\[
  \mathrm{X}_{k+1}^\bullet\longrightarrow\mathrm{X}_k^\bullet
\]
compatível com as filtrações e que faz de
$\bb{X}^\bullet=(\mathrm{X}_k^\bullet)_{k\in\bb{Z}}$ um sistema
projetivo de complexos.
\oldpage[0\textsubscript{III}]{76}
Podemos supor além disso que este sistema projetivo é estrito.
Com efeito, por \sref{0.13.7.1.1},
\[
  A_k
  \simeq
  A_{k+1}/\mathrm{F}^{k+1}(A_{k+1}).
\]
Assim, ao construir $\mathrm{X}_{k+1}^\bullet$, podemos tomar a
resolução injetiva de
$A_{k+1}/\mathrm{F}^{k+1}(A_{k+1})$ \emph{igual} a
$\mathrm{X}_k^\bullet$.
Segue-se então de (M, V, 2.3) que o morfismo de complexos
$\mathrm{X}_{k+1}^\bullet\longrightarrow\mathrm{X}_k^\bullet$ pode
ser construído de modo que, ao passar aos quocientes, induza um
isomorfismo
\[
  \mathrm{X}_{k+1}^\bullet/
  \mathrm{F}^{k+1}(\mathrm{X}_{k+1}^\bullet)
  \simto
  \mathrm{X}_k^\bullet
\]
que respeita as filtrações; esta é precisamente a condição de
\sref{0.13.5.1}.
\end{env}

\begin{env}[13.7.3]
\label{0.13.7.3}
Por construção, o sistema projetivo
$\mathrm{T}(\bb{X}^\bullet)$ de complexos
$\mathrm{T}(\mathrm{X}_k^\bullet)$ satisfaz as hipóteses de
\sref{0.13.5.3}.
Os resultados de \sref{0.13.5.4}, \sref{0.13.5.5} e
\sref{0.13.5.6} aplicam-se, portanto, às sucessões espectrais
\[
  \mathrm{E}\bigl(\mathrm{T}(\mathrm{X}_k^\bullet)\bigr)
  =
  \mathrm{E}(A_k).
\]
Para $1\leq r\leq+\infty$, escrevemos
$\mathrm{E}_r^{pq}(\bb{A})$ em lugar de
$\mathrm{E}_r^{pq}(\mathrm{T}(\bb{X}^\bullet))$
(cf.\ \sref{0.13.5.7} quando $r=+\infty$), e utilizamos a mesma convenção
para a notação análoga.
Em particular,
\begin{equation}
\label{0.13.7.3.1}
\tag{13.7.3.1}
  \mathrm{E}_1^{pq}(\bb{A})
  =
  \RR^{p+q}\mathrm{T}\bigl(\mathrm{gr}^p(\bb{A})\bigr),
\end{equation}
por \sref{0.13.6.4.2} e o fato de que o sistema projetivo
$(\mathrm{gr}^p(A_k))_{k\in\bb{Z}}$ é essencialmente constante.
Estes resultados, junto com \sref{0.13.4.3}, dão a proposição
seguinte, demonstrada pela primeira vez por Shih Weishu mediante um método
diferente (inédito).
\end{env}

\begin{proposition}[13.7.4]
\label{0.13.7.4}
\emph{(Shih).}
Seja $n$ um inteiro.
As duas condições seguintes são equivalentes:
\begin{enumerate}
  \item[{\rm a)}] Para todo par $(p,q)$ tal que $p+q=n$, o
  sistema projetivo
  $(\mathrm{E}_r^{pq}(\bb{A}))_{r\geq2}$ é essencialmente constante.
  \item[{\rm b)}] O sistema projetivo
  \[
    \RR^n\mathrm{T}(\bb{A})
    =
    \bigl(\RR^n\mathrm{T}(A_k)\bigr)_{k\in\bb{Z}}
  \]
  satisfaz {\rm(ML)}.
\end{enumerate}
Além disso, quando essas condições são satisfeitas, existe um
isomorfismo canônico
\begin{equation}
\label{0.13.7.4.1}
\tag{13.7.4.1}
  \mathrm{gr}^p\bigl(\RR^n\mathrm{T}(\bb{A})\bigr)
  \isoto
  \mathrm{E}_\infty^{p,n-p}(\bb{A})
  \qquad(p\in\bb{Z}).
\end{equation}
\end{proposition}

\begin{proof}
Por \sref{0.13.5.6}, a condição {\rm a)} diz precisamente que, para
$p+q=n$, o sistema projetivo
$(\mathrm{E}_\infty^{pq}(A_k))_{k\in\bb{Z}}$ é essencialmente
constante.
Por outro lado,
$\mathrm{gr}^p(\RR^n\mathrm{T}(A_k))$ é canonicamente isomorfo a
$\mathrm{E}_\infty^{p,n-p}(A_k)$.
Segue-se de \sref{0.13.4.3} que {\rm a)} e {\rm b)} são
equivalentes, e \sref{0.13.7.4.1} não é senão
\sref{0.13.5.6.1} no caso considerado.
\end{proof}

\begin{corollary}[13.7.5]
\label{0.13.7.5}
Seja $\sh{F}=(\sh{F}_k)_{k\in\bb{N}}$ um sistema projetivo de
feixes de grupos abelianos que satisfaz as condições {\rm(i)}, {\rm(ii)}
e {\rm(iii)} de \sref{0.13.3.1}, e ponhamos
$\sh{F}=\varprojlim_k\sh{F}_k$.
Suponhamos que, para o funtor
$\sh{G}\longmapsto\Gamma(X,\sh{G})$, o sistema projetivo
$(\mathrm{E}_r^{pq}(\sh{F}))_r$ é essencialmente constante sempre que
$p+q=n$ ou $p+q=n+1$.
Sobre $\HH^{n+1}(X,\sh{F})$, consideremos a filtração
\[
  \mathrm{F}^p\bigl(\HH^{n+1}(X,\sh{F})\bigr)
  =
  \Ker\bigl(
    \HH^{n+1}(X,\sh{F})
    \longrightarrow
    \HH^{n+1}(X,\sh{F}_{p-1})
  \bigr).
\]
Então existe um isomorfismo canônico
\begin{equation}
\label{0.13.7.5.1}
\tag{13.7.5.1}
  \mathrm{gr}^p\bigl(\HH^{n+1}(X,\sh{F})\bigr)
  \isoto
  \mathrm{E}_\infty^{p,n-p+1}(\sh{F})
  \qquad(p\in\bb{Z}).
\end{equation}
\end{corollary}

\begin{proof}
\oldpage[0\textsubscript{III}]{77}
Apliquemos \sref{0.13.7.4} tomando como $\cat{C}$ a categoria de feixes de
grupos abelianos sobre $X$, como $\cat{C}'$ a categoria de grupos abelianos,
e $\mathrm{T}=\Gamma$.
Dá, para todo $p\in\bb{Z}$, um isomorfismo canônico
\[
  \mathrm{gr}^p\bigl(\RR^{n+1}\Gamma(\sh{F})\bigr)
  \isoto
  \mathrm{E}_\infty^{p,n-p+1}(\sh{F}).
\]
Além disso, por \sref{0.13.7.4}, o sistema projetivo
$(\HH^n(X,\sh{F}_k))_{k\in\bb{Z}}$ satisfaz {\rm(ML)}; portanto
\sref{0.13.3.1} dá um isomorfismo canônico
\[
  \HH^{n+1}(X,\sh{F})
  \isoto
  \varprojlim_k\HH^{n+1}(X,\sh{F}_k).
\]

Como o sistema projetivo
$\RR^{n+1}\Gamma(\sh{F})$ satisfaz {\rm(ML)} por
\sref{0.13.7.4}, existe um isomorfismo canônico
\[
  \mathrm{gr}^p\bigl(\RR^{n+1}\Gamma(\sh{F})\bigr)
  \isoto
  \varprojlim_k
  \mathrm{gr}^p\bigl(\HH^{n+1}(X,\sh{F}_k)\bigr)
\]
por \sref{0.13.4.3}, e um isomorfismo canônico
\[
  \varprojlim_k
  \mathrm{gr}^p\bigl(\HH^{n+1}(X,\sh{F}_k)\bigr)
  \isoto
  \mathrm{gr}^p\left(
    \varprojlim_k\HH^{n+1}(X,\sh{F}_k)
  \right)
\]
por \sref{0.13.4.5}.
Resta apenas verificar que o isomorfismo precedente é
compatível com as filtrações de seus dois lados.
Isto segue-se imediatamente das definições e, para todo $p$,
da comutatividade de
\[
  \xymatrix{
    \HH^{n+1}(X,\sh{F})
      \ar[r]^-{\sim}\ar[dr]
    &
    \varprojlim_k\HH^{n+1}(X,\sh{F}_k)
      \ar[d]
    \\
    &
    \HH^{n+1}(X,\sh{F}_{p-1}).
  }
\]
\end{proof}

\begin{env}[13.7.6]
\label{0.13.7.6}
Seja $S$ um anel dotado de uma filtração
$(\mathrm{F}^i(S))_{i\in\bb{Z}}$ tal que
$\mathrm{F}^0(S)=S$ (e portanto
$\mathrm{gr}^i(S)=0$ para $i<0$).
Suponhamos que cada $A_k$, com a filtração definida em
\sref{0.13.7.1}, é um
$S$--$\cat{C}$-módulo filtrado no sentido de \sref{0.13.6.6}, e que os
morfismos $A_h\longrightarrow A_k$ para $k\leq h$ são morfismos de
$S$--$\cat{C}$-módulos filtrados.
Para abreviar, dizemos então que $\bb{A}$ é um
\emph{sistema projetivo de $S$--$\cat{C}$-módulos filtrados}.

É imediato que, para $k\leq h$ e
$1\leq r\leq+\infty$, os morfismos
\[
  \mathrm{B}_r^{pq}(A_h)\longrightarrow\mathrm{B}_r^{pq}(A_k),
  \qquad
  \mathrm{Z}_r^{pq}(A_h)\longrightarrow\mathrm{Z}_r^{pq}(A_k)
\]
são morfismos de
$\mathrm{gr}^\bullet(S)$--$\cat{C}'$-módulos bigraduados
(\sref{0.13.6.5}), e que, para $r$ finito, as famílias
\[
  \bigl(\mathrm{Z}_r^{pq}(\bb{A})\bigr),\qquad
  \bigl(\mathrm{B}_r^{pq}(\bb{A})\bigr),\qquad
  \bigl(\mathrm{E}_r^{pq}(\bb{A})\bigr)
\]
são $\mathrm{gr}^\bullet(S)$--$\cat{C}'$-módulos
bigraduados, sendo os dois primeiros submódulos de
$(\mathrm{E}_1^{pq}(\bb{A}))$.
Escrevemos de novo
$\mathrm{Z}_r^{(n)}(\bb{A})$,
$\mathrm{B}_r^{(n)}(\bb{A})$ e
$\mathrm{E}_r^{(n)}(\bb{A})$ para os respectivos subobjetos graduados
obtidos conservando unicamente os termos com $p+q=n$; estes são
$\mathrm{gr}^\bullet(S)$--$\cat{C}'$-módulos graduados.

Quando o sistema projetivo
$(\mathrm{E}_r^{pq}(\bb{A}))_r$ é essencialmente constante,
$\mathrm{E}_\infty^{pq}(\bb{A})$ é igualmente um
$\mathrm{gr}^\bullet(S)$--$\cat{C}'$-módulo bigraduado, e cada
$\mathrm{E}_\infty^{(n)}(\bb{A})$ é um
$\mathrm{gr}^\bullet(S)$--$\cat{C}'$-módulo graduado.
Além disso, para cada $k$, os isomorfismos
\[
  \beta^{p,n-p}:
  \mathrm{E}_\infty^{p,n-p}(A_k)
  \isoto
  \mathrm{gr}^p\bigl(\RR^n\mathrm{T}(A_k)\bigr)
\]
formam um isomorfismo de
$\mathrm{gr}^\bullet(S)$--$\cat{C}'$-módulos graduados de
$\mathrm{E}_\infty^{(n)}(A_k)$ sobre
$\mathrm{gr}^\bullet(\RR^n\mathrm{T}(A_k))$.
Sob as hipóteses precedentes, o isomorfismo
\[
  \beta^{p,n-p}:
  \mathrm{E}_\infty^{(n)}(\bb{A})
  \isoto
  \varprojlim_k
  \mathrm{gr}^\bullet\bigl(\RR^n\mathrm{T}(A_k)\bigr)
\]
respeita, portanto, também estas estruturas.
O mesmo vale manifestamente para o isomorfismo canônico
\[
  \mathrm{gr}^\bullet\bigl(\RR^n\mathrm{T}(\bb{A})\bigr)
  \isoto
  \varprojlim_k
  \mathrm{gr}^\bullet\bigl(\RR^n\mathrm{T}(A_k)\bigr);
\]
por conseguinte, os isomorfismos \sref{0.13.7.4.1} formam um
isomorfismo de
$\mathrm{gr}^\bullet(S)$--$\cat{C}'$-módulos graduados.
\end{env}

\begin{proposition}[13.7.7]
\label{0.13.7.7}
Seja $S$ um anel noetheriano $\mathfrak{J}$-ádico.
Suponhamos que $\cat{C}$ é uma categoria abeliana na qual todo objeto
é isomorfo a um subobjeto de um objeto injetivo, e seja
$\mathrm{T}$ um funtor covariante aditivo de $\cat{C}$ na
categoria de grupos abelianos.
Seja $\bb{A}=(A_k)_{k\in\bb{Z}}$ um
\oldpage[0\textsubscript{III}]{78}
sistema projetivo estrito de $S$--$\cat{C}$-módulos filtrados, para a
filtração $\mathfrak{J}$-ádica sobre $S$, limitado inferiormente.
Para um inteiro fixo $n$, suponhamos que seja satisfeita a condição seguinte:
\begin{itemize}
  \item[$(\mathrm{F}_n)$]
  \emph{O $\mathrm{gr}^\bullet(S)$-módulo graduado}
  \[
    \mathrm{E}_1^{(m)}(\bb{A})
    =
    \bigl(
      \RR^m\mathrm{T}(\mathrm{gr}^p(\bb{A}))
    \bigr)_{p\in\bb{Z}}
  \]
  \emph{de \sref{0.13.7.3.1} é de tipo finito para
  $m=n$ e $m=n+1$.}
\end{itemize}
Então:
\begin{enumerate}
  \item[{\rm(i)}] Os sistemas projetivos
  $(\RR^n\mathrm{T}(A_k))_{k\in\bb{Z}}$ e
  $(\RR^{n+1}\mathrm{T}(A_k))_{k\in\bb{Z}}$ satisfazem {\rm(ML)}.
  \item[{\rm(ii)}] Se
  \[
    \RR'^{\,n}\mathrm{T}(\bb{A})
    =
    \varprojlim_k\RR^n\mathrm{T}(A_k),
  \]
  então $\RR'^{\,n}\mathrm{T}(\bb{A})$ é um $S$-módulo de tipo
  finito.
  \item[{\rm(iii)}] A filtração sobre
  $\RR'^{\,n}\mathrm{T}(\bb{A})$ definida por
  \[
    \mathrm{F}^p\bigl(\RR'^{\,n}\mathrm{T}(\bb{A})\bigr)
    =
    \Ker\bigl(
      \RR'^{\,n}\mathrm{T}(\bb{A})
      \longrightarrow
      \RR^n\mathrm{T}(A_{p-1})
    \bigr)
    \qquad(p\in\bb{Z})
  \]
  é $\mathfrak{J}$-boa; isto é,
  \[
    \mathfrak{J}\,
    \mathrm{F}^p\bigl(\RR'^{\,n}\mathrm{T}(\bb{A})\bigr)
    \subset
    \mathrm{F}^{p+1}\bigl(\RR'^{\,n}\mathrm{T}(\bb{A})\bigr)
  \]
  para todo $p$, com igualdade para todo $p$ suficientemente grande.
  Em particular, a topologia sobre
  $\RR'^{\,n}\mathrm{T}(\bb{A})$ definida por esta filtração é a
  topologia $\mathfrak{J}$-ádica.
  \item[{\rm(iv)}] O sistema projetivo
  $(\mathrm{E}_r^{pq}(\bb{A}))_r$ é essencialmente constante sempre que
  $p+q=n$ ou $p+q=n+1$.
  Assim, $\mathrm{E}_\infty^{pq}(\bb{A})$ está definido
  (\sref{0.13.5.7}), e os isomorfismos
  \begin{equation}
  \label{0.13.7.7.1}
  \tag{13.7.7.1}
    \mathrm{gr}^p\bigl(\RR'^{\,n}\mathrm{T}(\bb{A})\bigr)
    \isoto
    \mathrm{E}_\infty^{p,n-p}(\bb{A})
    \qquad(p\in\bb{Z})
  \end{equation}
  formam um isomorfismo de
  $\mathrm{gr}^\bullet(S)$-módulos graduados.
\end{enumerate}
\end{proposition}

\begin{proof}
Em vista de \sref{0.13.7.7.1}, e levando em conta os
isomorfismos \sref{0.13.7.4.1}, por abuso de notação escreveremos também
$\RR^n\mathrm{T}(\bb{A})$ para o limite projetivo
$\RR'^{\,n}\mathrm{T}(\bb{A})$ do sistema projetivo
$(\RR^n\mathrm{T}(A_k))$.

O anel graduado $\mathrm{gr}^\bullet(S)$ é noetheriano
(Bourbaki, \emph{Alg.\ comm.}, cap.~III, \textsection2,
n\textsuperscript{o}~9, cor.~5 do teorema~2).
Portanto, a sucessão crescente de
$\mathrm{gr}^\bullet(S)$-submódulos graduados
$\mathrm{B}_r^{(m)}(\bb{A})$ de
$\mathrm{E}_1^{(m)}(\bb{A})$ é estacionária para $m=n$ e
$m=n+1$.
A condição {\rm(b)} de \sref{0.11.1.10} fica, portanto, satisfeita.
Segue-se que a condição {\rm(a)} de \sref{0.13.7.4} é satisfeita para
$n$ e para $n+1$, o que demonstra {\rm(i)}.

Além disso, os isomorfismos \sref{0.13.7.4.1}, junto com as
observações de \sref{0.13.7.6}, mostram que
$\mathrm{gr}^\bullet(\RR^n\mathrm{T}(\bb{A}))$ é um
$\mathrm{gr}^\bullet(S)$-módulo graduado isomorfo a
\[
  \mathrm{E}_\infty^{(n)}(\bb{A})
  =
  \mathrm{Z}_\infty^{(n)}(\bb{A})
  /
  \mathrm{B}_\infty^{(n)}(\bb{A}).
\]
Como $\mathrm{Z}_\infty^{(n)}(\bb{A})$ é um submódulo de
$\mathrm{E}_1^{(n)}(\bb{A})$, é de tipo finito, e também o é
$\mathrm{E}_\infty^{(n)}(\bb{A})$.
Finalmente, para a filtração
$(\mathrm{F}^p(\RR'^{\,n}\mathrm{T}(\bb{A})))$, segue-se de
\sref{0.13.4.5} que
\[
  \mathrm{gr}^\bullet\bigl(\RR^n\mathrm{T}(\bb{A})\bigr)
  \quad\hbox{e}\quad
  \mathrm{gr}^\bullet\bigl(\RR'^{\,n}\mathrm{T}(\bb{A})\bigr)
\]
são $\mathrm{gr}^\bullet(S)$-módulos graduados isomorfos.
Isto demonstra {\rm(iv)}.
As asserções {\rm(ii)} e {\rm(iii)} se seguirão destes resultados
e do lema seguinte.
\end{proof}

\begin{lemma}[13.7.7.2]
\label{0.13.7.7.2}
Seja $S$ um anel noetheriano $\mathfrak{J}$-ádico, e seja $M$
um $S$-módulo dotado de uma filtração codiscreta
$(\mathrm{F}^p(M))_{p\in\bb{Z}}$ tal que
\[
  \mathfrak{J}\mathrm{F}^p(M)
  \subset
  \mathrm{F}^{p+1}(M).
\]
Assim, $M$ é um módulo filtrado sobre $S$ dotado de seu
filtração $\mathfrak{J}$-ádica.
Suponhamos além disso que $M$ é separado para a topologia definida por
$(\mathrm{F}^p(M))$.
Então as condições seguintes são equivalentes:
\begin{enumerate}
  \item[{\rm a)}] $M$ é um $S$-módulo de tipo finito e
  $(\mathrm{F}^p(M))$ é uma filtração $\mathfrak{J}$-boa.
  \item[{\rm b)}] $\mathrm{gr}^\bullet(M)$ é um
  $\mathrm{gr}^\bullet(S)$-módulo de tipo finito.
  \item[{\rm c)}] Os $\mathrm{gr}^p(M)$ são $S$-módulos de
  tipo finito e, para todo $p$ suficientemente grande, os homomorfismos
  canônicos
  \begin{equation}
  \label{0.13.7.7.3}
  \tag{13.7.7.3}
    \mathfrak{J}\otimes_S\mathrm{gr}^p(M)
    \longrightarrow
    \mathrm{gr}^{p+1}(M)
  \end{equation}
  \oldpage[0\textsubscript{III}]{79}
  são sobrejetivos.
  Estão induzidos por
  $\mathfrak{J}\otimes_S\mathrm{F}^p(M)
  \longrightarrow\mathrm{F}^{p+1}(M)$, levando em conta o fato
  de que a imagem do homomorfismo composto
  \[
    \mathfrak{J}\otimes_S\mathrm{F}^{p+1}(M)
    \longrightarrow
    \mathfrak{J}\otimes_S\mathrm{F}^p(M)
    \longrightarrow
    \mathrm{F}^{p+1}(M)
  \]
  é
  $\mathfrak{J}\mathrm{F}^{p+1}(M)
  \subset\mathrm{F}^{p+2}(M)$.
\end{enumerate}

Para a demonstração, veja-se Bourbaki, \emph{Alg.\ comm.}, cap.~III,
\textsection3, n\textsuperscript{o}~1, prop.~3.
\end{lemma}

\begin{env}[13.7.7.4]
\label{0.13.7.7.4}
Para aplicar o lema \sref{0.13.7.7.2}, resta observar que a
topologia sobre $\RR'^{\,n}\mathrm{T}(\bb{A})$ definida pela filtração
considerada o torna um $S$-módulo separado e completo:
esta é a topologia de limite projetivo dos grupos discretos
$\RR^n\mathrm{T}(A_k)$.
Por outro lado, se $A_k=0$ para $k<k_0$, então também
$\RR^n\mathrm{T}(A_k)=0$ para $k<k_0$; portanto
\[
  \mathrm{F}^{k_0}\bigl(\RR'^{\,n}\mathrm{T}(\bb{A})\bigr)
  =
  \RR'^{\,n}\mathrm{T}(\bb{A}).
\]
Assim, todas as hipóteses do lema ficam efetivamente satisfeitas.
\end{env}

\begin{corollary}[13.7.8]
\label{0.13.7.8}
Se a hipótese $(\mathrm{F}_n)$ é satisfeita, então para todo $f\in S$
existe um isomorfismo canônico
\begin{equation}
\label{0.13.7.8.1}
\tag{13.7.8.1}
  \varprojlim_k\bigl(\RR^n\mathrm{T}(A_k)\bigr)_f
  \isoto
  \RR^n\mathrm{T}(\bb{A})\otimes_S S_{\{f\}}.
\end{equation}
\end{corollary}

\begin{proof}
Com efeito, $\RR^n\mathrm{T}(\bb{A})$ é um $S$-módulo de tipo finito,
e $S_{\{f\}}$ é uma $S$-álgebra ádica noetheriana
(\sref[0\textsubscript{I}]{0.7.6.11}), a saber, o completamento separado
de $S_f$ para a topologia $\mathfrak{J}$-pré-ádica
(\sref[0\textsubscript{I}]{0.7.6.2}).
Segue-se de \sref[0\textsubscript{I}]{0.7.7.8} e
\sref[0\textsubscript{I}]{0.7.7.1} que
\[
  \RR^n\mathrm{T}(\bb{A})\otimes_S S_{\{f\}}
\]
é isomorfo ao completamento separado de
$\RR^n\mathrm{T}(\bb{A})\otimes_S S_f$ para a
topologia $\mathfrak{J}$-pré-ádica.
Um sistema fundamental de vizinhanças de zero para esta topologia
é
\[
  \bigl(
    \mathfrak{J}^p\RR^n\mathrm{T}(\bb{A})
  \bigr)\otimes_S S_f;
\]
os grupos
\[
  \mathrm{F}^p\bigl(\RR^n\mathrm{T}(\bb{A})\bigr)
  \otimes_S S_f
\]
formam, portanto, também um sistema tal.
Este último grupo é o núcleo da aplicação canônica
\[
  \bigl(\RR^n\mathrm{T}(\bb{A})\bigr)_f
  \longrightarrow
  \bigl(\RR^n\mathrm{T}(A_{p-1})\bigr)_f.
\]
Por conseguinte, o grupo separado associado a
$\RR^n\mathrm{T}(\bb{A})\otimes_S S_f$ identifica-se com um subgrupo
$G$ de
\[
  \varprojlim_k\bigl(\RR^n\mathrm{T}(A_k)\bigr)_f.
\]
Mas o sistema projetivo
$((\RR^n\mathrm{T}(A_k))_f)_k$ satisfaz manifestamente {\rm(ML)}, e a
imagem de $(\RR^n\mathrm{T}(\bb{A}))_f$ em cada
$(\RR^n\mathrm{T}(A_k))_f$ é igual à imagem comum dos
$(\RR^n\mathrm{T}(A_h))_f$ para todo $h\geq k$ suficientemente grande.
Segue-se imediatamente que $G$ é \emph{denso em toda parte} em
$\varprojlim_k(\RR^n\mathrm{T}(A_k))_f$.
Como este último grupo é separado e completo, fica demonstrado o
corolário.
\end{proof}
